\chapter{Resultados}
\label{chap:resultados}
Neste capítulo são apresentados e analisados os resultados obtidos a partir dos experimentos descritos no Capítulo~\ref{chap:metodos}. A análise é organizada em quatro etapas progressivas: (i) validação da qualidade das correspondências automáticas entre os datasets\acrshort{ctmar} e DeepLesion/Stroke EIT, assegurando a confiabilidade da base de dados utilizada; (ii) estudos de ablation em treinos curtos para comparação arquitetural dos módulos de atenção e contexto multi-escala; (iii) seleção objetiva do modelo candidato com base em critérios de consistência e significância estatística; e (iv) avaliação do modelo selecionado em treinos de duração intermédia e prolongada.

Foi realizado um estudo para verificar qual o melhor peso de máscara para a função de perda. Foram testados os pesos 0.1, 0.05, e 0.01. A finalidade era observar quais os limites de atuação da segunda saída em regularização e preservação de detalhes. O modelo com peso 0.05 para a máscara apresentou melhor desempenho em validação, com ganhos marginais em\acrshort{psnr} e\acrshort{ssim}, sugerindo que a supervisão adicional contribui para uma correção mais equilibrada dos artefatos. O modelo com peso 0.1 apresentou resultados piores que o modelo \textit{Single-Head}, demonstrando que a supervisão excessiva da máscara pode levar à degradação de detalhes ao competir com a tarefa de reconstrução. Dessa forma, o peso da função de perda da saída de máscara foi diminuído para 0.05. Para esse peso, os resultados melhoram para\acrshort{ms-ssim},\acrshort{ssim},\acrshort{psnr} e\acrfull{mae} em relação ao peso 0.1, mas ainda apresentam uma leve degradação em relação ao modelo \textit{Single-Head}. O valor do DICE teve uma leve queda, mas manteve-se acima de 0.70, o que é considerado um desempenho aceitável para a tarefa de segmentação de artefatos com finalidade de regularização. Para o peso 0.01, o desempenho melhora em relação ao modelo com peso 0.05, mais uma vez melhorando\acrshort{ms-ssim},\acrshort{ssim},\acrshort{psnr} e\acrshort{mae}, mas o valor do\Acrfull{dice} tem uma queda mais acentuada, ficando abaixo de 0.70.

